\hypertarget{control-flow-preprocessor}{%
\chapter{CONTROL FLOW \& PREPROCESSOR}\label{control-flow-preprocessor}}

\hypertarget{advanced-control-flow}{%
\chapter{ADVANCED CONTROL FLOW}\label{advanced-control-flow}}

\hypertarget{ternary-operator}{%
\section{9.1 Ternary Operator}\label{ternary-operator}}

\begin{lstlisting}[language={C++}]
#include <iostream>
using namespace std;

int main() {
    int x = 10, y = 5;
    
    // condition ? true_value : false_value
    int max = (x > y) ? x : y;
    cout << "Max: " << max << endl;  // 10
    
    // Nested ternary (use with caution)
    int age = 20;
    string status = (age < 18) ? "Minor" : (age < 65) ? "Adult" : "Senior";
    cout << status << endl;
    
    // String ternary
    cout << (x % 2 == 0 ? "Even" : "Odd") << endl;
    
    return 0;
}
\end{lstlisting}

\hypertarget{goto-statement-avoid}{%
\section{9.2 goto Statement (Avoid)}\label{goto-statement-avoid}}

\begin{lstlisting}[language={C++}]
#include <iostream>
using namespace std;

int main() {
    // goto is generally discouraged
    int x = 0;
    
loop:
    cout << x << " ";
    x++;
    
    if (x < 5) {
        goto loop;
    }
    cout << endl;
    
    // Better alternative: use loops
    for (int i = 0; i < 5; i++) {
        cout << i << " ";
    }
    cout << endl;
    
    return 0;
}
\end{lstlisting}

\hypertarget{label-goto-for-error-handling}{%
\section{9.3 Label \& Goto for Error
Handling}\label{label-goto-for-error-handling}}

\begin{lstlisting}[language={C++}]
#include <iostream>
#include <cstdlib>
using namespace std;

int main() {
    FILE* file = NULL;
    char* buffer = NULL;
    
    // Using goto for cleanup (rare acceptable use)
    file = fopen("test.txt", "r");
    if (!file) {
        cout << "Failed to open file" << endl;
        goto cleanup;
    }
    
    buffer = new char[100];
    if (!buffer) {
        cout << "Memory allocation failed" << endl;
        goto cleanup;
    }
    
    // Do work...
    
cleanup:
    if (buffer) delete[] buffer;
    if (file) fclose(file);
    
    return 0;
}
\end{lstlisting}

\begin{center}\rule{0.5\linewidth}{0.5pt}\end{center}

\hypertarget{preprocessor-directives}{%
\chapter{PREPROCESSOR DIRECTIVES}\label{preprocessor-directives}}

\hypertarget{define-and-include}{%
\section{7.1 \#define and \#include}\label{define-and-include}}

\begin{lstlisting}[language={C++}]
// Define constants
#define PI 3.14159
#define MAX_SIZE 100
#define SQUARE(x) ((x) * (x))

// Conditional compilation
#define DEBUG

#ifdef DEBUG
    #define LOG(msg) cout << msg << endl
#else
    #define LOG(msg)  // Do nothing in release
#endif

#include <iostream>
using namespace std;

int main() {
    cout << "PI = " << PI << endl;
    
    int arr[MAX_SIZE];
    cout << "Array size: " << sizeof(arr) << endl;
    
    cout << "Square of 5: " << SQUARE(5) << endl;
    
    LOG("Debug message");
    
    return 0;
}
\end{lstlisting}

\hypertarget{conditional-compilation}{%
\section{7.2 Conditional Compilation}\label{conditional-compilation}}

\begin{lstlisting}[language={C++}]
#include <iostream>
using namespace std;

// Platform-specific code
#ifdef _WIN32
    #define OS "Windows"
#elif __APPLE__
    #define OS "macOS"
#elif __linux__
    #define OS "Linux"
#else
    #define OS "Unknown"
#endif

int main() {
    cout << "Running on: " << OS << endl;
    
#if defined(DEBUG)
    cout << "Debug mode" << endl;
#else
    cout << "Release mode" << endl;
#endif
    
    return 0;
}
\end{lstlisting}

\hypertarget{pragma-directives}{%
\section{7.3 Pragma Directives}\label{pragma-directives}}

\begin{lstlisting}[language={C++}]
#include <iostream>
using namespace std;

// Disable specific warnings
#pragma warning(disable: 4996)  // MSVC

// Pack structure
#pragma pack(1)
struct PackedData {
    char a;
    int b;
    double c;
};
#pragma pack()

int main() {
    cout << "Size of PackedData: " << sizeof(PackedData) << endl;
    // Without pragma pack: 24 (aligned)
    // With pragma pack: 13 (packed)
    
    return 0;
}
\end{lstlisting}

\begin{center}\rule{0.5\linewidth}{0.5pt}\end{center}

\hypertarget{inline-functions-macros}{%
\chapter{INLINE FUNCTIONS \& MACROS}\label{inline-functions-macros}}

\hypertarget{macro-functions-vs-inline-functions}{%
\section{12.1 Macro Functions vs Inline
Functions}\label{macro-functions-vs-inline-functions}}

\begin{lstlisting}[language={C++}]
#include <iostream>
using namespace std;

// Macro function - preprocessor substitution
#define ADD_MACRO(a, b) ((a) + (b))

// Inline function - type-safe
inline int add_inline(int a, int b) {
    return a + b;
}

int main() {
    cout << ADD_MACRO(5, 3) << endl;         // 8
    cout << add_inline(5, 3) << endl;       // 8
    
    // Macro danger: side effects
    int x = 5, y = 3;
    cout << ADD_MACRO(x++, y++) << endl;    // Evaluates as: ((x++) + (y++))
    cout << "x = " << x << ", y = " << y << endl;  // x = 6, y = 4
    
    // Inline function is safer
    x = 5, y = 3;
    cout << add_inline(x++, y++) << endl;   // 8
    cout << "x = " << x << ", y = " << y << endl;  // x = 6, y = 4 (correct)
    
    return 0;
}
\end{lstlisting}

\begin{center}\rule{0.5\linewidth}{0.5pt}\end{center}
